\documentclass[12pt]{article}

%----------------- Packages ------------------
\usepackage[utf8]{inputenc}
\usepackage[T1]{fontenc}
\usepackage{amsmath, amssymb}
\usepackage{geometry}
\usepackage{xcolor}
\usepackage{titlesec}
\usepackage{fancyhdr}
\usepackage[breakable]{tcolorbox}
\usepackage{graphicx}
\usepackage{tikz}
\usepackage{float}
\usetikzlibrary{arrows.meta, decorations.markings}

%----------------- Page Setup -----------------
\geometry{a4paper, margin=2.5cm}
\pagestyle{fancy}
\fancyhf{}
\rhead{Final Exam — \textbf{2019}}
\lhead{Analysis IV}
\cfoot{\thepage}

%----------------- Title Styling --------------
\titleformat{\section}{\normalfont\Large\bfseries}{Exercise \thesection:}{1em}{}
\titleformat{\subsection}{\normalfont\bfseries}{Answer:}{1em}{}

%----------------- Custom Boxes ----------------
\tcbuselibrary{listingsutf8}
\newtcolorbox{answerbox}{
  colback=gray!10,
  colframe=black,
  fonttitle=\bfseries,
  title=Answer Area,
  breakable,
  before skip=10pt,
  after skip=10pt
}

%----------------- Document Start --------------
\begin{document}

%----------------- Exam Info -------------------
\begin{center}
  \Large\textbf{UNIVERSITY NAME} \\[1em]
  \large\textit{Analysis IV — Final Exam} \\[0.5em]
  \large\textit{Year: 2019} \\[2em]
\end{center}

\vspace{0.5cm}

%----------------- EXERCISE 1 ------------------
\section{}
Determine the nature (convergence/divergence) of the following sequences:

\[
\begin{aligned}
&a) \; u_n = \cos \left( \pi \frac{2n^3 + n^2 + a n + b}{2n^2} \right), \quad a, b \in \mathbb{R} \\
&b) \; u_n = \left( 1 - \frac{1}{n^a} \right)^n, \quad a > 0
\end{aligned}
\]

\begin{answerbox}

\end{answerbox}

%----------------- EXERCISE 2 ------------------
\section{}
Consider the series of functions \( \sum_{n \geq 1} f_n \) where

\[
f_n(x) = \frac{x}{n(1+n^2 x)}.
\]

Let \( f(x) = \sum_{n=1}^{\infty} f_n(x) \) when it exists.

\begin{enumerate}
    \item Show that $f$ is defined and continious on  \( \mathbb{R}^+ \setminus \{0\} \).
    \item Show that there exists a unique \( \lambda > 0 \) such that \( \lim_{x \to \infty}f(x) = \lambda \).
    \item Show that $f$ is differentiable on \( \mathbb{R}^+_*\) and compute \( f'(x) \). 
    \item Sketch the graph of \( f \).
\end{enumerate}

\begin{answerbox}

\end{answerbox}

%----------------- EXERCISE 3 ------------------
\section{}
Define the sequence \( (u_n) \) by \( u_0 = u_1 = 1 \) and for \( n \in \mathbb{N} \):

\[
u_{n+2} = u_{n+1} + 2u_n + (-1)^n.
\]

Let \( f(x) = \sum_{n=0}^{\infty} u_n x^n \) when it exists, and let \( R \) be its radius of convergence.

\begin{enumerate}
    \item Show that \( \forall n \in \mathbb{N} \),
    \[
    1 \leq u_n \leq 2^{n+1} - 1
    \]
    and deduce that \( R \geq \frac{1}{2} \).
    \item Show that for \( |x| < \min(1, R)\), \( f(x) \) satisfies:
    \[
    f(x) = \frac{x^2 + x + 1}{(x+1)^2 (1 - 2x)}.
    \]
    \item Give an explicit expression for \( u_n \) as a function of \( n \).
\end{enumerate}

\begin{answerbox}

\end{answerbox}

%----------------- END ------------------
\end{document}